\markboth{DISCUSSION}{DISCUSSION}
\section{Discussion}\label{sec:discussion}

This section discusses architectural trade-offs, the scope boundaries we deliberately set, the limitations of the current system, and the practical lessons we drew from building Training Copilot.

\subsection{Architectural Trade-offs}\label{sec:discussion:tradeoffs}

\textbf{Multi-agent vs.\ single-agent design.} The multi-agent architecture introduces coordination overhead: each query traverses the orchestrator, one or more specialists, and their MCP servers, accumulating latency from multiple LLM calls and network round-trips, where a single-agent design would eliminate the A2A layer entirely. Three benefits justified the added complexity. With over 40 tools across all servers (before the optional 116-tool Telegram bridge), a single agent frequently selected incorrect tools---research confirms larger tool sets increase misselection~\citep{qu_tool_2024, patil_gorilla_2023}, and scoping each specialist to at most three servers substantially improved accuracy. Each specialist also carries a domain-optimised prompt (the recovery specialist interprets HRV relative to baseline, the route specialist geocodes before routing); a single system prompt embedding all domain knowledge was unwieldy and the model often failed to follow it, which the modular prompt architecture~\citep{wang_survey_2024} solved. Cross-domain queries also benefit from parallel execution, since the orchestrator's concurrent \texttt{ask\_*} calls reduce latency to the slowest specialist rather than the sum.

\textbf{MCP vs.\ direct API integration.} Encapsulating each API behind an MCP server introduces a process boundary---every tool call is an HTTP round-trip---that direct integration would eliminate, but the abstraction provides three properties that justify the cost, consistent with established microservice principles~\citep{dragoni_microservices_2017}: independence (each server manages its own authentication, rate limiting, and retry logic, so restarting one does not affect others), uniform discovery (adding the Telegram proxy and flythrough servers required one file and one configuration line each, confirming the single-line-extensibility promise), and deployment flexibility (servers are independently containerisable, with URLs swapped via environment variables for Docker Compose).

\subsection{Scope and Focus}\label{sec:discussion:scope}

Following the midterm review, we deliberately shifted from broadening the feature surface to deepening the capabilities at the heart of the product. Most subsequent effort therefore went not into new integrations but into hardening the two things we regard as Training Copilot's core: \textit{agentic data analytics}---turning raw, multi-source telemetry into synthesised, grounded insight---and \textit{agentic training and recovery planning}, understood in both senses of recovery: guidance for returning to training after injury, and day-to-day readiness inferred from physiological signals such as HRV, sleep, and Body Battery. This is where our evaluation concentrates its emphasis (Section~\ref{sec:evaluation}): the personas and scorers are built to stress the agentic reasoning path rather than the presentation layer. The remaining surfaces---the analytics dashboards, the Telegram delivery channel, and the 3D flythrough---are intentionally lightweight add-ons, shaped to our own preferences as endurance athletes and straightforward to retune or replace; they demonstrate the architecture's extensibility rather than constitute the contribution, which is why their rough edges (Table~\ref{tab:known_issues}) are acceptable at this stage.

Training Copilot is a research prototype validating two hypotheses---that MCP-based data integration achieves single-line extensibility, and that multi-agent coordination produces contextually richer recommendations than monolithic approaches. Concretely, the path we prioritised is the \textit{core conversational analysis and planning workflow}, from a natural-language question through agent coordination to a data-grounded, synthesised answer, while accepting known incompleteness in peripheral features. The capabilities we consider production-quality for a research prototype are training-readiness assessment (combining recovery, weather, and calendar data), training-load analysis (CTL/ATL/TSB trends), route planning with interactive maps, RAG-grounded fitness advice with source citations, and schedule-aware planning with calendar event creation. Several features are architecturally present but not fully polished; Table~\ref{tab:known_issues} documents them as conscious scope decisions rather than oversights.

\begin{table}[ht]
\begin{center}
\begin{minipage}{\textwidth}
\caption[Known incomplete areas]{Known incomplete areas (placeholder, pending a final pre-submission re-assessment).}\label{tab:known_issues}
\begin{tabularx}{\textwidth}{X}
\toprule
\textbf{Known incomplete areas} \\
\midrule
\textit{To be finalised before submission.} The rough edges noted at the midterm (inconsistent chart triggering, chat-map metric overlays, within-session context windowing, relative-date handling for calendar writes, and sport-aware route timing) have since been addressed, and no new ones surfaced in this pass. Any that remain at submission will be re-assessed and listed here. \\
\botrule
\end{tabularx}
\end{minipage}
\end{center}
\end{table}


\subsection{Limitations}\label{sec:discussion:limitations}

\textbf{Latency.} Complex cross-domain queries typically take 8--15 seconds end-to-end in our development setup---acceptable for training planning but not for in-workout guidance. LLM inference time dominates; MCP round-trips contribute under 500\,ms in aggregate. Token-level streaming would improve perceived responsiveness but is currently disabled for compatibility with the KIT university gateway, which has shown instabilities under streaming load.

\textbf{Knowledge currency.} The RAG-based fitness specialist draws from fifty public-domain books from Project Gutenberg, necessarily historical. While progressive overload, periodisation, and recovery principles are timeless~\citep{grgic_periodization_2017}, modern exercise science has advanced substantially, and contemporary open-access publications would improve the specialist's relevance.

\textbf{Platform dependencies.} Training Copilot depends on external APIs of varying availability. Strava now requires a subscription for Standard Tier API access~\citep{strava_api_update_2026}, and Garmin Connect has no official public API, so the third-party library used may break with platform changes; the MCP architecture mitigates this, since each server is independently replaceable, but does not eliminate the dependency.

\textbf{Route quality and data availability.} The most popular consumer route planners---Komoot and Outdooractive---keep their curated, community-contributed route catalogues behind closed doors: that route data \textit{is} their product, so declining to expose it through a public API is a rational commercial choice rather than an oversight. Training Copilot therefore plans over OpenRouteService, an open engine on raw OpenStreetMap data; it can synthesise a loop of any target distance from any start point, even where no pre-planned route exists (Section~\ref{sec:data:other}), but it optimises over the road and path graph rather than serving hand-picked scenic itineraries, so a suggestion's scenic quality can trail a Komoot route. This is a limitation of the available data ecosystem, not of our integration, which would accept a richer route provider through the same one-server interface the moment one opens up. Because the project's focus is agentic analytics and training/recovery planning, the search for a higher-quality route source is paused for now and left to future work.

\textbf{Single-user focus.} Multi-tenant deployment with per-user credential vaults and session isolation is architecturally supported but not yet implemented. The security implications---SSRF prevention, tool-output sanitisation, egress control for user-added MCP servers, encrypted token storage---remain open; the MCP-specific threat surface that \cite{hou_mcp_2025} catalogue (tool poisoning, prompt injection, and lifecycle attacks on third-party servers) applies directly once users can register their own servers, and addressing it is prerequisite to any public deployment.

\subsection{Ethical Considerations}\label{sec:discussion:ethics}

An AI system that provides training recommendations bears responsibility for athlete safety. Training Copilot explicitly does not provide medical advice; system prompts instruct all specialists to recommend professional consultation for injury-related questions or abnormal metrics. Consumer wearables have known accuracy limitations~\citep{fuller_reliability_2020}, mitigated partially through multi-source corroboration (cross-referencing HRV, sleep, and Body Battery rather than any single metric) but not eliminated. The fitness RAG corpus also reflects early 20th-century biases and limited demographic scope; expanding it with contemporary, demographically inclusive literature is a priority for future work.

\subsection{Lessons Learned}\label{sec:discussion:lessons}

\textbf{Tool docstrings are the interface.} A tool's docstring quality has a disproportionate impact on selection accuracy---vague descriptions led to incorrect calls, while precise descriptions stating \textit{when} to call a tool substantially improved accuracy, consistent with \cite{patil_gorilla_2023}'s findings on API documentation quality. We treated docstrings as the primary interface between model and data layer.

\textbf{Graceful degradation is not optional.} With seven external dependencies, partial outages are the norm. The three-layer strategy---skip missing MCP servers, report unavailable specialists, surface tool errors as structured data---proved essential when individual services restarted or returned partial data during development; what we initially treated as a nice-to-have quickly became load-bearing.

\textbf{Record full, clip for the model.} The recording-and-clipping pattern (Section~\ref{sec:agents:recording}) was one of our most impactful design decisions. GPS streams and intraday timelines can exceed 100\,KB; rendering the UI from the full trace rather than the model's response preserves data fidelity without burdening the LLM---a pattern we recommend for any tool-augmented agent handling large structured results.

\textbf{Scoping constrains and enables.} Restricting each specialist to its own MCP servers began as a simplification but proved a design strength: it prevents tool confusion, reduces prompt complexity, and makes each specialist independently testable.

\textbf{Source preservation requires explicit engineering.} The orchestrator's synthesis step dropped citations in roughly one-third of fitness-related answers until we added an explicit \texttt{ensure\_sources()} post-processing step---automated synthesis and citation preservation are fundamentally in tension, and bridging that gap takes deliberate engineering.
